%%!TEX root = the-kleiner.tex
\chapter{Barycenters}\label{chap:barycenter}

Barycenters provide yet another way to use your Euclidean intuition for $\CAT(0)$ spaces.
We will apply them to dimension theory.


\section{Definition}

Let us denote by $\triangle^k\subset \RR^{k+1}$\index{1@$\triangle^m$ (simplex)}
the \index{standard simplex}\emph{standard $k$-simplex}; 
that is, $\bm{\mu}=(\mu_0,\ldots,\mu_k)\in\triangle^k$ if $\mu_0+\ldots+\mu_k=1$ and $\mu_i\ge0$ for all $i$.

Consider a point array $\bm{p}=(p_0,\ldots,p_k)$ in the Euclidean space $\EE^n$.
Recall that
\[z=\mu_0\cdot p_0+\ldots+\mu_k\cdot p_k\]
is called the \index{barycenter}\emph{barycenter} of the point array $\bm{p}=(p_0,\ldots,p_k)$ with masses $\bm{\mu}\z=(\mu_0,\ldots, \mu_k)\in \triangle^k$.
Equivalently, 
\[z\df \argmin (\mu_0\cdot f_0+\ldots+\mu_k\cdot f_k),\eqlbl{eq:f-spx}\]
where $f_i=\tfrac12\cdot\distfun_{p_i}^2$ for each $i$, and $\argmin f$\index{3@$\argmin$ (point of minimum)} denotes a point of minimum of the function $f$.

The map $\spx{}\:\triangle^k\mapsto \EE^n$ defined by $\spx{}\:\bm{\mu}\mapsto z$ is called the \index{barycentric simplex} \emph{barycentric simplex} of the array $\bm{p}$.
If needed we may denote this map by $\spx{\bm{p}}$ or, more generally, $\spx{\bm{f}}$.
The latter means that we define the map via \ref{eq:f-spx} for an array of functions $\bm{f}=(f_0,\ldots,f_k)$.
Formally speaking, the definition \ref{eq:f-spx} makes sense for any array of functions in a metric space;
but in this case, the map might be undefined or non-uniquely defined.

Further, we will work with this definition in $\CAT(0)$ spaces.
We will essentially use only that on a geodesic $\CAT(0)$ space
functions of the type $f=\tfrac12\cdot\distfun_{p}^2$ are 1-convex; see \ref{ex:convex-distfun}.

\section{Barycentric simplex}

\begin{thm}{Theorem}\label{thm:barycenter}
Let $\spc{X}$ be a complete geodesic space 
and $\bm{f}\z=(f_0,\z\ldots,f_k)\:\spc{X}\to\RR^{k+1}$
be an array of non-negative 1-convex locally Lipschitz functions.
Then the barycentric simplex 
$\spx{\bm{f}}\:\triangle^k\to \spc{X}$ is a uniquely defined Lipschitz map.

In particular, we have that for any point array $\bm{p}\z=(p_0,\z\ldots,p_k)$ in a complete geodesic $\CAT(0)$ space
the barycentric simplex $\spx{\bm{p}}$ is a uniquely defined Lipschitz map.
\end{thm}


\begin{thm}{Lemma}\label{lem:argmin(convex)}
Suppose $\spc{X}$ is a complete geodesic space and $f\:\spc{X}\z\to\RR$ is a locally Lipschitz, $1$-convex function.
Then $\argmin f$ is uniquely defined.
\end{thm}

\parit{Proof.}
Note that
\begin{clm}{}\label{midpoint}
 if $z$ is a midpoint of a geodesic $[x y]$, then 
\[s\le f(z)
\le
\tfrac{1}{2}\cdot f(x)+\tfrac{1}{2}\cdot f(y)-\tfrac{1}{8}\cdot\dist[2]{x}{y}{},
\]
where $s$ is the infimum of $f$.
\end{clm}

\parit{Uniqueness.}
Assume that $x$ and $y$ are distinct minimum points of $f$.
Let $z$ be the midpoint of a geodesic $[x y]$.
From \ref{midpoint} we have
\[f(z)<f(x)=f(y)\] 
--- a contradiction. 

\parit{Existence.}
Fix a point $p\in \spc{X}$, and
let $\Lip\in\RR$ be a Lipschitz constant of $f$ in a neighborhood of $p$.

Choose a sequence of points $p_n\in \spc{X}$ such that $f(p_n)\to s$.
Applying \ref{midpoint} for $x\z=p_n$, $y\z=p_m$, we see that $p_n$ is a Cauchy sequence.
Thus the sequence $p_n$ converges to a minimum point of $f$.
\qeds

Assume $\phi$ is a convex function defined on a real interval $\II$.
For $t_0\in\II$, let us define the \index{right derivative}\emph{right} (respectively, \index{left derivative}\emph{left}) \index{3@$\phi^\pm$ (right/left derivative)}\emph{derivative} $\phi^+(t_0)$ ($\phi^-(t_0)$) at $t_0$ by
\[\phi^\pm(t_0)=\lim_{t\to t_0\pm} \frac{\phi(t)-\phi(t_0)}{|t-t_0|}.\]
Note that our sign convention for $\phi^-$ is not standard --- for $\phi(t)=t$ we have
$\phi^+(t)=1$ and $\phi^-(t)=-1$.

\parit{Proof of \ref{thm:barycenter}.}
Since each $f_i$ is $1$-convex, for any $\bm{\lambda}=(\lambda_0,\ldots,\lambda_k)\z\in\triangle^k$
the convex combination 
\[\left(\sum_i \lambda_i\cdot f_i\right)\:\spc{X}\to\RR\]
is also $1$-convex.
Therefore, according to \ref{lem:argmin(convex)}, the barycentric simplex 
%$\spx{\bm{f}}(\bm{\lambda})$ is defined for any $\bm{\lambda}\in\triangle^k$.
$\spx{\bm{f}}$ is uniquely defined on $\triangle^k$.
 
For $\bm{\lambda},\bm{\mu}\in\triangle^k$,
let 
\begin{align*}
f_{\bm{\lambda}}
&=\sum_i \lambda_i\cdot f_i,
&
f_{\bm{\mu}}
&=\sum_i \mu_i\cdot f_i,
\\
p
&=\spx{\bm{f}}(\bm{\lambda}),
&
q
&=\spx{\bm{f}}(\bm{\mu}),
\end{align*}
Choose a geodesic $\gamma$ from $p$ to $q$;
suppose $s=\dist{p}{q}{}$ and so $\gamma(0)=p$ and $\gamma(s)=q$.
Observe the following:
\begin{itemize}
\item The function $\phi(t)=f_{\bm{\lambda}}\circ\gamma(t)$ has a minimum at $0$.
Therefore $\phi^+(0)\ge 0$.

\item The function $\psi(t)=f_{\bm{\mu}}\circ\gamma(t)$ has a minimum at $s$.
Therefore $\psi^-(s)\ge 0$.
\end{itemize}
From $1$-convexity of $f_{\bm{\mu}}$, we have
\[\psi^+(0)+\psi^-(s)+s\le0.\]

Let $\Lip$ be a Lipschitz constant for all $f_i$ in a neighborhood $\Omega\ni p$.
Then 
\[\psi^+(0)
\le 
\phi^+(0)+\Lip\cdot\|\bm{\lambda}-\bm{\mu}\|_1,\]
where $\|\bm{\lambda}-\bm{\mu}\|_1=\sum_{i=0}^k|\lambda_i-\mu_i|$.
It follows that given $\bm{\lambda}\in\triangle^k$, there is a constant $\Lip$ such that
\begin{align*}
\dist{\spx{\bm{f}}(\bm{\lambda})}{\spx{\bm{f}}(\bm{\mu})}{}
&=
s
\le
\\
&\le 
\Lip\cdot\|\bm{\lambda}-\bm{\mu}\|_1
\end{align*}
for any $\bm{\mu}\in\triangle^k$.
In particular, there is $\eps>0$ such that if $\|\bm{\lambda}-\bm{\mu}\|_1<\eps,$ $\|\bm{\lambda}-\bm{\nu}\|_1 <\eps$, then $\spx{\bm{f}}(\bm{\mu})$, $\spx{\bm{f}}(\bm{\nu})\in\Omega$.
Thus the same argument as above implies 
\[\dist{\spx{\bm{f}}(\bm{\mu})}{\spx{\bm{f}}(\bm{\nu})}{}
\le \Lip\cdot\|\bm{\mu}-\bm{\nu}\|_1\]
for any $\bm{\mu}$ and $\bm{\nu}$ sufficiently close to $\bm{\lambda}$; that is, $\spx{\bm{f}}$ is locally Lipschitz.
Since $\triangle^k$ is compact, $\spx{\bm{f}}$ is Lipschitz.
\qeds

\begin{thm}{Exercise}\label{ex:finite-action-CAT}
Let $G$ be a subgroup of the group of isometries of a proper geodesic $\CAT(0)$ space.
Assume that
\begin{subthm}{ex:finite-action-CAT:finite}
$G$ is finite, or, more generally, 
\end{subthm}
\begin{subthm}{ex:finite-action-CAT:compact}
 $G$ is compact.
\end{subthm}
Show that the action of $G$ has a fixed point.
\end{thm}

\begin{thm}{Advanced exercise}\label{ex:bary-jensen}
Let $z$ be the barycenter of points $p_0,\ldots,p_k$ equipped with masses $\mu_0,\ldots,\mu_k$; we assume that $\mu_i\ge0$ for all $i$,
$\mu_0+\ldots+\mu_k=1$, and the points $p_0,\ldots,p_k$ are in a complete geodesic $\CAT(0)$ space $\spc{U}$.

\begin{subthm}{ex:bary-jensen:moment}
Show that
\[\sum_i \mu_i\cdot \dist[2]{q}{p_i}{}\ge \dist[2]{z}{q}{} +\sum_i \mu_i\cdot \dist[2]{z}{p_i}{}\]
for any $q\in \spc{U}$.
\end{subthm}

\begin{subthm}{ex:bary-jensen:moment+}
 Show that
\[\sum_i
\mu_i\cdot \dist[2]{z}{p_i}{}
\le
\sum_{i<j}
\mu_i\cdot\mu_j\cdot \dist[2]{p_i}{p_j}{}.\]
\end{subthm}


\begin{subthm}{ex:bary-jensen:jensen}
(Jensen inequality.)
Try to show that
\[h(z)\le \mu_0\cdot h(p_0)+\ldots+\mu_k\cdot h(p_k)\]
for any convex function $h\:\spc{U}\to\RR$.
\end{subthm}

\end{thm}


\section{Convexity of up-set}

\begin{thm}{Definition}\label{def:ordung}
For two real arrays $\bm{v}$, $\bm{w}\in \RR^{k+1}$,
$\bm{v}=(v_0,\z\ldots,v_k)$
and 
$\bm{w}=(w_0,\z\ldots,w_k)$,
we will write
$\bm{v}\succcurlyeq\bm{w}$ if $v_i\ge w_i$ for each $i$.
\end{thm}

Given a subset $Q\subset \RR^{k+1}$,
denote by $\Up Q$ \label{PAGE.def:Up}
the smallest upper set containing $Q$;
that is,
\begin{align*}
\Up Q 
&=
\set{\bm{v}\in\RR^{k+1}}{\exists\, \bm{w}\in Q\ \text{such that}\ \bm{v}\succcurlyeq\bm{w}},
\end{align*}

\begin{thm}{Proposition}\label{thm:up-convex}
Let $\spc{X}$ be a complete geodesic space 
and $\bm{f}\z=(f_0,\z\ldots,f_k)\:\spc{X}\to\RR^{k+1}$
be an array of non-negative 1-convex locally Lipschitz functions.
Consider the set $W=\Up[\bm{f}(\spc{X})]\subset\RR^{k+1}$.
Then 


\begin{subthm}{thm:up-convex:convex}
The set $W$ is convex.
\end{subthm}

\begin{subthm}{thm:up-convex:bry}
$\bm{f}[\spx{\bm{f}}(\triangle^k)]\subset\partial W$.
Moreover, $\bm{f}[\spx{\bm{f}}(\triangle^k)\setminus \spx{\bm{f}}(\partial\triangle^k)]$ is an open set in $\partial W$.
\end{subthm}

\begin{subthm}{thm:up-convex:bry+}
$W=\Up(\bm{f}[\spx{\bm{f}}(\triangle^k)])$; 
in other words, $\Up(\bm{f}[\spx{\bm{f}}(\triangle^k)])\supset\bm{f}(\spc{X})$.
In particular, $W$ is closed.
\end{subthm}

\end{thm}


\begin{figure}[t!]
\vskip0mm
\centering
\includegraphics{mppics/pic-1106}
\caption{}
\end{figure}

\parit{Proof.}
Let $V=\bm{f}(\spc{X})\subset\RR^{k+1}$; so $W\z=\Up V$.
Denote by $\bar V$ the closure of $V$.

\parit{\ref{SHORT.thm:up-convex:convex}.}
Convexity of all $f_i$ implies that
for any two points $p,q\in \spc{X}$ and $t\z\in[0,1]$ we have
\[(1-t)\cdot\bm{f}(p)+t\cdot \bm{f}(q)
\succcurlyeq
\bm{f}\circ\gamma(t),\]
where $\gamma$ denotes a geodesic path from $p$ to $q$. 
Therefore, $W$ is convex.

\parit{\ref{SHORT.thm:up-convex:bry}+\ref{SHORT.thm:up-convex:bry+}.}
Choose $p\in \spx{\bm{f}}(\triangle^k)$.
Note that if $\bm{f}(p)\succcurlyeq\bm{w}$ for some $\bm{w}\in W$, then $\bm{f}(p)\succcurlyeq\bm{w}\succcurlyeq\bm{f}(x)$ for some $x\in \spc{X}$ and hence $x=p$ by the uniqueness of $\spx{\bm{f}}$.
Therefore, $\bm{f}(p)=\bm{w}$.
It follows that $\bm{f}(p)\in\partial W$;
therefore $\bm{f}[\spx{\bm{f}}(\triangle^k)]$ lies in a convex hypersurface $\partial W$.

Choose $\bm{w}\in W$.
Observe that $\bm{w}\succcurlyeq\bm{v}$ for some $\bm{v}\in \bar V\cap \partial W$.
Note that $W$ is supported at $\bm{v}$ by a hyperplane 
\[\Pi=\set{(x_0,\ldots,x_k)\in \RR^{k+1}}{\mu_0\cdot x_0+\ldots+\mu_k\cdot x_k=\const}\]
for some $\bm{\mu}=(\mu_0,\ldots,\mu_k)\in\triangle^{k}$.
Let $p=\spx{\bm f}(\bm{\mu})$.
By \ref{lem:argmin(convex)}, $\bm{f}(p)\z=\bm{v}$;
in particular $\bm{v}\in V$.

Note that $p\in \spx{\bm{f}}(\triangle^k)\setminus\spx{\bm{f}}(\partial\triangle^k)$ if and only if 
$\bm{f}(p)$ is supported by a plane as above for some $\bm{\mu}\in\triangle^{k}$,
but not supported by any such plane with $\bm{\mu}\in\partial\triangle^{k}$.
This condition is open, therefore $\spx{\bm{f}}(\triangle^k)\z\setminus\spx{\bm{f}}(\partial\triangle^k)$ is an open set.

Since $\triangle^k$ is compact, $W=\Up(\bm{f}[\spx{\bm{f}}(\triangle^k)])$ is closed.
\qeds

\section{Nondegenerate simplex}

Given an array $\bm{f}=(f_0,\ldots,f_k)$,
we denote by $\bm{f}^{\without i}$ the subarray of $\bm{f}$ with $f_i$ removed;
that is, 
\[\bm{f}^{\without i\,}
\df
(f_0,\ldots,f_{i-1},f_{i+1},\ldots,f_k).\]
Note that $\spx{\bm{f}^{\without i}}$ coincides with the restriction of $\spx{\bm{f}}$ to the corresponding facet of $\triangle^k$.

If $\Im \spx{\bm{f}}$ is not covered by the union of $\Im \spx{\bm{f}^{\without i}}$ over all $i$, then we say that $\spx{\bm{f}}$ is nondegenerate.
In other words, $\spx{\bm{f}}$ is \index{nondegenerate simplex}\emph{nondegenerate} if 
\[\spx{\bm{f}}(\triangle^k)\setminus \spx{\bm{f}}(\partial \triangle^k)\ne\emptyset.\]

\begin{thm}{Exercise}\label{ex:barysimple}
Let $\spc{U}$ be a complete geodesic $\CAT(0)$ space.

Show that the image of the 1-dimensional barycentric simplex for a pair of points $p_0,p_1\in\spc{U}$ is the geodesic $[p_0\,p_1]$.

Construct a $\CAT(0)$ space with a three-point array $(p_0,p_1,p_2)$ such that its barycentric simplex is nondegenerate and noninjective.
\end{thm}

\begin{thm}[!]{Exercise}\label{lem:nondeg-test-with-balls}
Let $\bm{p}\z=(p_0,\ldots, p_k)$ be a point array
in a complete length $\CAT(0)$ space $\spc{U}$, 
and $B_i=\cBall[p_i,r_i]$ for some array of positive reals $(r_0,r_1,\ldots,r_k)$.

\begin{subthm}{lem:nondeg-test-with-ball:U}
Suppose $\bigcap_i B_i\ne \emptyset$.
Show that 
\[\Im\spx{\bm{p}}\subset \bigcup_i B_i.\]
\end{subthm}

\begin{subthm}{lem:nondeg-test-with-balls:nondeg}
Suppose $\bigcap_i B_i= \emptyset$, but $\bigcap_{i\ne j} B_i\ne \emptyset$ for any $j$.
Show that 
$\spx{\bm{p}}$ is nondegenerate.
\end{subthm}

\begin{subthm}{lem:nondeg-test-with-balls:nondeg+}
Suppose $\spx{\bm{p}}$ is nondegenerate.
Show that the condition in \ref{SHORT.lem:nondeg-test-with-balls:nondeg} holds for some
 array of positive reals $(r_0,\ldots,r_k)$.
\end{subthm}

\end{thm}



\section{bi-H\"older embedding}

\begin{thm}{Theorem}\label{thm:bihoelder}
Let $\spc{X}$ be a complete geodesic space 
and $\bm{f}\z=(f_0,\z\ldots,f_k)\:\spc{X}\to\RR^{k+1}$
be an array of 1-convex locally Lipschitz functions.
Then the set 
\[Z=\spx{\bm{f}}(\triangle^k)\setminus \spx{\bm{f}}(\partial \triangle^k)\]
is $C^{\frac12}$-bi-H\"older to an open domain in $\RR^k$.
\end{thm}

\parit{Proof.}
Let $\proj\:\RR^{k+1}\to\Pi$ be orthogonal projection to the hyperplane $x_0+\ldots+x_k=0$.
We will show that the restriction $\proj\circ \bm{f}|_Z$ is a bi-H\"older embedding.

Since the map $\proj\circ \bm{f}$ is Lipschitz,
it is sufficient to construct its right inverse and show that it is $C^{\frac12}$-continuous.

Given $\bm{v}=(v_0,\ldots,v_k)\in\Pi$, consider the function
$h_{\bm{v}}\: \spc{X}\to \RR$ defined by
\[h_{\bm{v}}(p)=\max_i\{f_i(p)-v_i\}.\]
Note that $h_{\bm{v}}$ is $1$-convex.
Let 
$$\map(\bm{v})\df\argmin h_{\bm{v}}.$$
According to Lemma~\ref{lem:argmin(convex)}, $\map(\bm{v})$ is uniquely defined.

If $\bm{v}=\proj \bm{f}(p)$, then 
\[f_i\circ \map(\bm{v})\le f_i(p)\]
for any $i$.
In particular, if $p\in \spx{\bm{f}}(\triangle^k)$, then $p=\map(\bm{v})$.
That is, $\map$ is a right inverse of the restriction $\bm{f}|_{\spx{\bm{f}}(\triangle^k)}$.

Given $\bm{v},\bm{w}\in\RR^{k+1}$,
set $p=\map (\bm{v})$ and $q=\map (\bm{w})$.
Since $h_{\bm{v}}$ and $h_{\bm{w}}$ are 1-convex, we have
\begin{align*}
h_{\bm{v}}(q)
&\ge 
h_{\bm{v}}(p)+\tfrac{1}{2}\cdot\dist[2]{p}{q}{},
&
h_{\bm{w}}(p)
&\ge 
h_{\bm{w}}(q)+\tfrac{1}{2}\cdot\dist[2]{p}{q}{}.
\end{align*}
Therefore,
\begin{align*}
\dist[2]{p}{q}{}
&\le 
2\cdot\sup_{x\in\spc{X}}\{ |h_{\bm{v}}(x)-h_{\bm{w}}(x)| \}
\le
\\
&\le 
2\cdot\max_{i}\{|v_i-w_i|\}.
\end{align*}
In particular,
$\map$ is $C^{\frac{1}{2}}$-continuous.

Finally, by \ref{thm:up-convex:bry}, $\bm{f}(Z)$ is a $k$-dimensional manifold --- hence the result.
\qeds

\section{Review of topological dimension}

Let $\spc{X}$ be a metric space and $\set{V_\beta}{\beta\in\mathcal{B}}$
 be an open cover of $\spc{X}$.
Let us recall two notions in general topology:
\begin{itemize}

\item The \index{order of a cover}\emph{order} of $\set{V_\beta}{\beta\in\mathcal{B}}$ is the supremum of all integers $n$ such that there is a collection of $n+1$ elements of $\{V_\beta\}$ with nonempty intersection.

\item An open cover $\set{W_\alpha}{\alpha\in\IndexSet}$ of $\spc{X}$ is called a \index{refinement of a cover}\emph{refinement} of $\set{V_\beta}{\beta\in\mathcal{B}}$ if for any $\alpha\in\IndexSet$ there is $\beta\in\mathcal{B}$ such that $W_\alpha\subset V_\beta$.

\end{itemize}

\begin{thm}{Definition}\label{def:TopDim}\index{dimension!topological dimension}\index{topological dimension}
Let $\spc{X}$ be a metric space. 
The topological dimension of $\spc{X}$ is defined to be the minimum of non-negative integers $n$
such that for any open cover of $\spc{X}$ there is a locally finite open refinement with order~$n$.
If no such $n$ exists, then the topological dimension of $\spc{X}$ is infinite.

The topological dimension of $\spc{X}$ will be denoted by $\TopDim\spc{X}$.
\end{thm}

The invariants satisfying statements \ref{dim-axiom-norm} and \ref{dim-axiom-sigma} are commonly called the \index{dimension}\emph{dimension};
for that reason, we call these statements axioms.

\begin{thm}{Normalization axiom}
\label{dim-axiom-norm} For any $m\in\ZZ_{\ge0}$,
\[\TopDim\EE^m=m.\]

\end{thm}

\begin{thm}{Cover axiom}\label{dim-axiom-sigma} 
If $\{A_n\}_{n=1}^\infty$ is a countable closed cover of $\spc{X}$, then
\begin{align*}
\TopDim \spc{X}&=\sup\nolimits_n\{\TopDim A_n\}.
\end{align*}

\end{thm}

\parbf{On product spaces.} 
The following inequality holds for arbitrary metric spaces
\begin{align*}
\TopDim (\spc{X}\times\spc{Y})
&\le 
\TopDim \spc{X}+ \TopDim\spc{Y}
\end{align*}
and this inequality might be strict \cite{pontyagin-surface}.

\medskip

\begin{thm}{Definition}
Let $\spc{X}$ be a metric space
and $F\:\spc{X}\to\RR^m$ be a continuous map.
A point $\bm{z}\in \Im F$ is called a \emph{stable value} of $F$
if there is $\eps>0$ such that $\bm{z}\in\Im F'$ 
for any \emph{$\eps$-close} to $F$ continuous map $F'\:\spc{X}\to\RR^m$,
that is, $|F'(x)-F(x)|<\eps$ for all $x\in \spc{X}$.
\end{thm}



The next theorem follows from \cite[theorems VI 1$\&$2]{hurewicz-wallman}.
(This theorem also holds for non-separable metric spaces \cite{nagata}, \cite[3.2.10]{engelking}). 

{\sloppy

\begin{thm}{Stable value theorem}\label{thm:stable-value}
Let $\spc{X}$ be a separable metric space.
Then $\TopDim\spc{X}\ge m$ if and only if there is a map $F\:\spc{X}\to\RR^m$ with a stable value.
\end{thm}

}

\section{Dimension theorem}\label{sec:dim-CAT}

\begin{thm}{Theorem}\label{thm:dim-infty-CBA}
For any proper geodesic $\CAT(0)$ space $\spc{U}$, the following statements are equivalent:

\begin{subthm}{thm:dim-infty-CBA:TopDim}
\[\TopDim \spc{U}\ge m.\]
\end{subthm}

\begin{subthm}{thm:dim-infty-CBA:bary} 
For some $z\in \spc{U}$ there is an array of $m+1$ balls $B_i\z=\oBall(a_i,r_i)$
such that 
\[\bigcap_i B_i=\emptyset,
\quad\text{but}\quad
\bigcap_{i\ne j} B_i\ne \emptyset
\quad \text{for each $j$}.\]

\end{subthm}


\begin{subthm}{thm:dim-infty-CBA:mnfld} 
There is a $C^{\frac{1}{2}}$-embedding of an open set in $\RR^m$ into $\spc{U}$;
that is, $\map$ is bi-Hölder with exponent $\tfrac{1}{2}$.
\end{subthm}

\end{thm}


\begin{thm}{Lemma}\label{lem:approximation-cba}
Let $\spc{U}$ be a proper geodesic $\CAT(0)$ space
and $\rho\:\spc{U}\z\to\RR$ be a continuous positive function.
Then there is a locally finite countable simplicial complex $\spc{N}$,
a locally Lipschitz map $\map\:\spc{U}\z\to \spc{N}$, 
and a Lipschitz map $\Psi\:\spc{N}\to\spc{U}$ such that:

\begin{subthm}{lem:approximation-cba:displacement}
The displacement of the composition $\Psi\circ\map\:\spc{U}\to\spc{U}$ is bounded by $\rho$;
that is,
\[\dist{x}{\Psi\circ\map(x)}{}<\rho(x)\] 
for any $x\in\spc{U}$.
\end{subthm}

\begin{subthm}{lem:approximation-cba:im}
If $\TopDim\spc{U}\le m$, 
then the $\Psi$-image of $\spc{N}$ 
coincides with the image of its $m$-skeleton.
\end{subthm}

\end{thm}

\parit{Proof.}
Choose a locally finite countable covering $\set{\Omega_\alpha}{\alpha\in\IndexSet}$ of $\spc{U}$ such that $\Omega_\alpha\subset \oBall(x,\tfrac{1}{3}\cdot\rho(x))$ for any $x\in \Omega_\alpha$. 

Denote by $\spc{N}$ the \index{nerve}\emph{nerve} of the covering $\{\Omega_\alpha\}$;
that is, $\spc{N}$ is an abstract simplicial complex with 
%set of vertices formed by $\IndexSet$,
vertex set $\IndexSet$,
such that a finite subset 
$\{\alpha_0,\z\ldots,\alpha_n\}\subset\IndexSet$
forms a simplex if and only if
\[\Omega_{\alpha_0}
\cap
\ldots\cap
\Omega_{\alpha_n}\ne\emptyset.\]

Choose a Lipschitz partition of unity 
$\phi_\alpha\:\spc{U}\to [0,1]$ subordinate to $\{\Omega_\alpha\}$.
Consider the map $\map\:\spc{U}\to \spc{N}$ such that the barycentric coordinate of $\map(p)$ is $\phi_\alpha(p)$.
Note that $\map$ is locally Lipschitz. 
Clearly, the $\map$-preimage of any open simplex in $\spc{N}$ lies in $\Omega_\alpha$ for some $\alpha\in\IndexSet$.

For each $\alpha\in\IndexSet$, 
choose $x_\alpha\in\Omega_\alpha$.
Let us extend the map $\alpha\mapsto x_\alpha$
to a map $\Psi\:\spc{N}\to\spc{U}$ that is barycentric on each simplex.
According to \ref{thm:barycenter}, this extension exists, 
and $\Psi$ is locally Lipschitz.

\parit{\ref{SHORT.lem:approximation-cba:displacement}.}
Fix $x\in\spc{U}$.
Denote by $\triangle$ the minimal simplex that contains $\map(x)$, 
and let $\alpha_0,\alpha_1,\ldots,\alpha_n$ be the vertices of $\triangle$.
%Denote by $\triangle$ the minimal simplex that contains $\map(x)$;
%and let $(\alpha_0,\alpha_1,\ldots,\alpha_n)$ be the vertices of $\triangle$.
Note that $\alpha$ is a vertex of $\triangle$ if and only if $\phi_{\alpha}(x)>0$.
Thus
\[\dist{x}{x_{\alpha_i}}{}<\tfrac{1}{3}\cdot\rho(x)\] 
for any $i$.
Therefore 
\[\diam\Psi(\triangle)
\le
\max_{i,j}\{\dist{x_{\alpha_i}}{x_{\alpha_j}}{}\}
<
\tfrac{2}{3}\cdot\rho(x).\]
In particular, 
\[\dist{x}{\Psi\circ\map(x)}{}\le\dist{x}{x_{\alpha_0}}{}+\diam \Psi(\triangle) <\rho(x).\]

\parit{\ref{SHORT.lem:approximation-cba:im}.}
Assume the contrary;
that is, $\Psi(\spc{N})$ is not contained in the $\Psi$-image of the $m$-skeleton of $\spc{N}$.
Then for some $k>m$,
there is a $k$-simplex $\triangle^k$ in $\spc{N}$
such that the barycentric simplex $\sigma=\Psi|_{\triangle^k}$ is nondegenerate;
that is, 
$$W=\Psi(\triangle^k)\setminus\Psi(\partial\triangle^k)\ne \emptyset.
$$
By \ref{thm:bihoelder}, $\TopDim\spc{U}\ge k$ --- a contradiction.
\qeds

\parit{Proof of \ref{thm:dim-infty-CBA}; \ref{SHORT.thm:dim-infty-CBA:bary}$\Rightarrow$\ref{SHORT.thm:dim-infty-CBA:mnfld}$\Rightarrow$\ref{SHORT.thm:dim-infty-CBA:TopDim}.}
The implication \ref{SHORT.thm:dim-infty-CBA:bary}$\Rightarrow$\ref{SHORT.thm:dim-infty-CBA:mnfld} follows from Exercise~\ref{lem:nondeg-test-with-balls}
and Theorem~\ref{thm:bihoelder}, and \ref{SHORT.thm:dim-infty-CBA:mnfld}$\Rightarrow$\ref{SHORT.thm:dim-infty-CBA:TopDim} is trivial.
 
\parit{\ref{SHORT.thm:dim-infty-CBA:TopDim}$\Rightarrow$\ref{SHORT.thm:dim-infty-CBA:bary}.}
According to \ref{thm:stable-value}, 
there is a continuous map $F\:\spc{U}\to \RR^m$ with a stable value.

Fix $\eps>0$.
Since $F$ is continuous, there is a continuous positive function $\rho$ defined on $\spc{U}$ such that 
\[\dist{x}{y}{}<\rho(x)
\quad\Rightarrow\quad
|F(x)- F(y)|<\tfrac13\cdot\eps.\]
Apply \ref{lem:approximation-cba} to $\rho$.
For the resulting simplicial complex $\spc{N}$ 
 and the maps $\map\:\spc{U}\z\to \spc{N}$, $\Psi\:\spc{N}\to \spc{U}$, we have
\[|F\circ \Psi\circ\map(x)-F(x)|<\tfrac13\cdot\eps\] 
for any $x\in \spc{U}$.

Arguing by contradiction,
assume $\TopDim\spc{U}<m$.
By \ref{lem:approximation-cba:im},
the image of $F\circ\Psi\circ\map$ lies in the $F\circ\Psi$-image of the $(m-1)$-skeleton of $\spc{N}$;
In particular, it can be covered by a countable collection of Lipschitz images of $(m-1)$-simplexes.
Hence
$\bm{0}\in \RR^m$ is not a stable value of $F\circ\Psi\circ\map$.
Since $\eps>0$ is arbitrary, we get the result.
\qeds

The following exercise is a generalization of Helly's theorem; for closely related statements see \cite[Prop. 5.3]{kleiner-1999} and \cite{ivanov2014}.


\begin{thm}{Exercise}\label{ex:helly}
Let $K_1,\ldots,K_n$ be closed convex subsets in a proper length $\CAT(0)$ space $\spc{U}$.
Suppose that $\TopDim \spc{U}\le m$ for some $m<n$ and any $m+1$ subsets from $\{K_1,\ldots, K_n\}$ have a common point.
Show that all subsets $K_1,\ldots,K_n$ have a common point.
\end{thm}

\section{Hausdorff dimension}

\begin{thm}{Definition}
\label{def:HausDim}\index{dimension!Hausdorff dimension}\index{Hausdorff dimension}
Let $\spc{X}$ be a metric space. 
Its \emph{Hausdorff dimension} is defined as
\[\HausDim\spc{X}=\sup\set{\alpha\in\RR}{\HausMes_\alpha(\spc{X})>0},\]
where $\HausMes_\alpha$ denotes the $\alpha$-dimensional Hausdorff measure.
\end{thm}

The following theorem follows from \cite[theorems V 8 and VII 2]{hurewicz-wallman}.

\begin{thm}{Szpilrajn's theorem}\label{thm:szpilrajn} 
Let $\spc{X}$ be a separable metric space.
Assume $\TopDim\spc{X}\ge m$.
Then $\HausMes_m \spc{X}>0$.

In particular, 
$\TopDim\spc{X}\le\HausDim\spc{X}$.
\end{thm}

{\sloppy

Except for Szpilrajn's theorem, there are no other relations between topological and Hausdorff dimension for general separable spaces.
Moreover, the following exercise implies that the same holds for compact geodesic $\CAT(0)$ spaces of topological dimension at least 1.

}

\begin{thm}{Exercise}\label{ex:dim-top-haus-CAT}
Given $\alpha\ge 1$ construct a metric on the binary tree such that it has compact completion of Hausdorff dimension $\alpha$.

Conclude that for any integer $m\ge 1$ and real $\alpha\ge m$ there is a compact $\CAT(0)$ space with topological dimension $m$ and Hausdorff dimension $\alpha$. 
\end{thm}


\section{Notes}

The dimension theorem is the main application of barycenters in Alexandrov geometry.
It is due to Bruce Kleiner \cite{kleiner-1999};
a more precise version was proved by Alexander Lytchak \cite{lytchak:diff}.
Kleiner conjectured that the dimension theorem holds for arbitrary complete geodesic $\CAT(\kappa)$ spaces \cite{kleiner-1999}; see also \cite[p.~133]{gromov-1993}.
For separable spaces, the answer is ``yes'', and it follows from Kleiner's argument \cite[Corollary 14.13]{alexander-kapovitch-petrunin-2025}.

One may wonder if bi-H\"older condition \ref{thm:dim-infty-CBA:mnfld} can be improved to bi-Lipschitz;
this is unknown even for compact spaces.
However if a compact geodesic $\CAT(0)$ space $\spc{U}$ has finite topological dimension $m$,
then a slight modification of Kleiner's technique can be used to show that there is a bi-Lipschitz embedding of an $m$-cube into $\spc{U}$ \cite[Theorem 14.15]{alexander-kapovitch-petrunin-2025}.
In particular, there is a bi-Lipschitz embedding of an $n$-cube for any $n\le m$.
If $\TopDim\spc{U}=\infty$, then we expect existence of a bi-Lipschitz embedding of an $n$-cube for any integer $n\ge1$.
The statement is trivial for $n=1$; in this case any geodesic gives an isometric embedding.
For $n=2$, it follows since minimal and metric minimizing surfaces in $\spc{U}$ are $\CAT(0)$;
see \cite{petrunin-stadler} and the references therein.
Any such surface is locally bi-Lipschitz to the Euclidean plane.
For $n\ge 3$ the question remains open.
% ??? Maybe it is solved by Stadler and Wenger --- we need to check later.

Giuliano Basso studied barycenters in spaces with a conical bicombing~\cite{basso}, which in particular include all geodesic $\CAT(0)$ spaces.

